% ECONOMICS_PROBLEM: {"id": "C73-651", "title": "Karlin's fictitious-play rate conjecture: disproved", "jel_code": "C73", "source_id": "OP-0146"}
% FIRST_POSTED: 2026-10-09
% ATTEMPT_STATUS: unresolved
% ENTRY_KIND: research_attempt
\documentclass[11pt]{article}
\usepackage[T1]{fontenc}
\usepackage[utf8]{inputenc}
\usepackage[margin=25mm]{geometry}
\usepackage{amsmath,amssymb,amsthm,mathtools,xurl,hyperref}
\newtheorem{proposition}{Proposition}
\newtheorem{lemma}{Lemma}
\newtheorem{theorem}{Theorem}
\newtheorem{corollary}{Corollary}
\newcommand{\R}{\mathbb R}
\newcommand{\E}{\mathbb E}
\newcommand{\N}{\mathbb N}
\newcommand{\ind}{\mathbf 1}
\DeclareMathOperator*{\argmax}{arg\,max}
\DeclareMathOperator*{\argmin}{arg\,min}
\setlength{\parindent}{0pt}
\setlength{\parskip}{6pt}
\title{Karlin's fictitious-play rate conjecture: disproved}
\author{GPT-6 (Codex)}
\date{9 October 2026}
\begin{document}
\maketitle
\section*{Historical status and scope of this attempt}
Karlin's universal square-root rate conjecture for ordinary fictitious
play is false. Daskalakis and Pan give identity-game counterexamples with
adversarial legal tie choices; Wang gives a later counterexample without
ties after the initial step. This attempt checks the exact rate
quantifiers, derives useful finite-time certificates and perturbation
bounds, and isolates why several attempted repairs do not prove the
historical claim. It does not independently reconstruct the full slow
trajectory in either paper, so its independent reconstruction remains
unresolved. That label does not reinstate a false historical conjecture
as an open problem.

For a maximizing row player and minimizing column player with matrix $A$,
write the empirical strategies after $t$ plays as $x^t,y^t$. The next
choices simultaneously best respond to those old empirical strategies,
and their empirical updates have denominator $t+1$. The gap is
\[
 g_A(x,y)=\max_i(Ay)_i-\min_j(x^TA)_j.
\]
Since $\min_j(x^TA)_j\leq x^TAy\leq\max_i(Ay)_i$, this gap is
nonnegative and bounds either player's unilateral payoff improvement.
Thus a small computed gap is a valid approximate-equilibrium certificate,
regardless of how the profile was generated.

\section*{Exact integer recurrences for the identity game}
Take $A=I_n$, and let $X(t)=tx^t$ and $Y(t)=ty^t$ be the integer action
counts. They are nonnegative and both sum to $t$. Simultaneous fictitious
play is exactly
\[
 \begin{aligned}
 i_{t+1}&\in\argmax_iY_i(t),& X(t+1)&=X(t)+e_{i_{t+1}},\\
 j_{t+1}&\in\argmin_jX_j(t),& Y(t+1)&=Y(t)+e_{j_{t+1}}.
 \end{aligned}                                                    \tag{1}
\]
Both decisions must use the old counts. Substituting the newly updated
$X(t+1)$ into the second decision would define a different process.
The gap has the exact form
\[
                 g_{I_n}(x^t,y^t)=
                 \frac{\max_iY_i(t)-\min_jX_j(t)}{t}.             \tag{2}
\]
In particular, an unnormalized count discrepancy of order $t^{1-1/n}$
corresponds to a normalized gap of order $t^{-1/n}$, not to a growing
equilibrium error.

For a transparent check, initialize the two-action identity game with
both players using action 1, and break later ties by the smallest index.
The first eight count states and gaps are
\[
 \begin{array}{c|c|c|c}
 t&X(t)&Y(t)&g(t)\\\hline
 1&(1,0)&(1,0)&1\\
 2&(2,0)&(1,1)&1/2\\
 3&(3,0)&(1,2)&2/3\\
 4&(3,1)&(1,3)&1/2\\
 5&(3,2)&(1,4)&2/5\\
 6&(3,3)&(1,5)&1/3\\
 7&(3,4)&(2,5)&2/7\\
 8&(3,5)&(3,5)&1/4
 \end{array}
\]
This verifies the recurrence and normalization, including the increase
from $t=2$ to $t=3$. It is neither a slow-rate counterexample nor evidence
for a universal asymptotic upper bound. In particular, a gap need not
decrease at every iteration even when it eventually converges to zero.

\section*{The precise asymptotic contradiction}
The strong proposed bound allows constants depending on a fixed game,
initialization, and legal tie-breaking trajectory, but requires
$g(t)\leq Ct^{-1/2}$ at every sufficiently large time. To refute it, one
fixed trajectory and an unbounded sequence of times are enough.

Suppose such a trajectory satisfies $g(t_k)\geq c t_k^{-\gamma}$,
where $c>0$, $t_k\to\infty$, and $\gamma<1/2$. Then
\[
                    \sqrt{t_k}\,g(t_k)
                    \geq c t_k^{1/2-\gamma}\longrightarrow\infty.
\]
No finite $C$ can satisfy the claimed eventual bound. Applied to the
published identity examples with $\gamma=1/n$, this contradiction uses
$n>2$; the exponent $1/2$ at $n=2$ alone gives no contradiction. Applied
to the later $10\times10$ example with exponent $1/3$, it also excludes
the conjectured rate when all best responses after initialization are
unique. Thus merely prescribing a better tie-breaking rule cannot repair
a universal theorem contradicted by the latter construction.

The lower bounds are asymptotic statements about specific trajectories.
They do not assert that fictitious play fails to converge, that every
trajectory in the exhibited game is slow, or that no restricted family
has the faster rate. Those stronger conclusions require different
quantifiers and are not inferred here.

\section*{Gap stability versus trajectory stability}
For any two same-sized payoff matrices with
$\|A-B\|_\infty=\max_{ij}|A_{ij}-B_{ij}|\leq\varepsilon$, and any
fixed probability vectors $x,y$, one has
\[
                   |g_A(x,y)-g_B(x,y)|\leq2\varepsilon.           \tag{3}
\]
Indeed every coordinate of $Ay-By$ has magnitude at most
$\varepsilon$, so their maxima differ by at most $\varepsilon$.
Similarly each coordinate of $x^TA-x^TB$ differs by at most
$\varepsilon$, and their minima obey the same bound. The triangle
inequality proves (3). The bound is independent of the number of actions.

It does not follow that two fictitious-play trajectories stay close.
Their chosen actions depend discontinuously on payoff comparisons at
ties. There is, however, a useful conditional finite-horizon result.
Fix a horizon $H\geq2$, an initial pure action pair, and a trajectory
of $A$ whose chosen
best responses at all decision times $1,\ldots,H-1$ are unique. Let
$\Delta_H>0$ be the minimum margin between each selected response and
its runner-up, over both players and those times, assuming both players
have at least two actions. If
$\|A-B\|_\infty<\Delta_H/2$, the trajectories agree through time $H$.

To prove this, induct on decision time. If the histories agree, so do
the current empirical strategies. Every pairwise payoff comparison
changes by at most $2\|A-B\|_\infty<\Delta_H$. Its strict sign remains
unchanged, so both selected responses agree and the next histories
agree. This completes the induction. A player with one action requires
no margin test. The theorem gives no uniform infinite-horizon stability
unless the margins have a positive lower bound independent of $H$.
In particular, perturbing a tied identity counterexample does not
by itself prove a tie-free asymptotic counterexample.

\section*{A profile certificate and the limitation of simulation}
Let $u=(1/n,\ldots,1/n)$. In the identity game put
$M=\max_jy_j$ and $m=\min_ix_i$. Since $M\geq1/n$ and $m\leq1/n$,
$g=M-m$ bounds both $M-1/n$ and $1/n-m$. The positive deviations of $y$
from $u$ have at most $n-1$ nonzero entries unless all vanish; the negative
deviations of $x$ satisfy the same count. Because deviations sum to zero,
\[
              \|y-u\|_1\leq2(n-1)g,
              \qquad\|x-u\|_1\leq2(n-1)g.                       \tag{4}
\]
Thus a computed small gap certifies proximity to the unique identity-game
equilibrium. This special error bound is not a general theorem equating
payoff residual with distance to equilibria in arbitrary games.

For every finite simulated horizon $H$, the number
$C_H=\max_{1\leq t\leq H}\sqrt t\,g(t)$ is finite. Consequently any
finite run is compatible with a square-root estimate using its own
constant $C_H$. The unresolved reconstruction step is to derive the
published infinite phase recurrence, prove all its best-response
inequalities, and show $C_H$ diverges along that recurrence. Recording
an increasing finite prefix alone does not establish that conclusion.

\section{Completion Estimate}
25\%

This is a post hoc, heuristic estimate for the full catalogue target.
The attempt checks identity-game recurrences and rate quantifiers and proves gap, finite-horizon trajectory and equilibrium-distance bounds. The historical conjecture is correctly reported as disproved, but an independent disproof still requires an infinite slow-trajectory construction with all best-response inequalities, including the tie-free qualification.

\section*{References}
Constantinos Daskalakis and Qinxuan Pan,
\emph{A Counter-Example to Karlin's Strong Conjecture for Fictitious Play},
FOCS 2014; arXiv:1412.4840.
\url{https://arxiv.org/abs/1412.4840}.
Yuanhao Wang, \emph{Tie-breaking Agnostic Lower Bound for Fictitious Play},
arXiv:2507.09902 (2025), Theorem 1 and the augmented-game construction.
\url{https://arxiv.org/pdf/2507.09902}.
These sources supply the historical counterexamples. The count audit,
finite-horizon stability proof, and profile bounds above state the
additional checks completed in this attempt.
\end{document}
